    Ab der 7. Programmzeile wird
\begin{code}
double n = 2*i + 1
n = 1.0/n;
\end{code}
  berechnet. Die Reihe wechselt ja ständig das Vorzeichen. Wenn wir nach $i=1000$ Schritten aufhören, dann ist der Fehler ungefähr so groß wie der nächste Term, den wir nicht mehr berechnet haben. Der nächste Term wäre
\begin{equation*}
  \Delta\pi = 4\frac{1}{2 n + 1} = \frac{4}{2001} \approx 0,002\,.
\end{equation*}
  Genau genommen ist der Fehler etwas kleiner, weil sich die restlichen Terme gegenseitig teilweise aufheben — aber für eine Schätzung reicht das.
